
\usepackage[margin=2.5cm]{geometry}
\usepackage{graphicx}
\usepackage{amsmath,amssymb}
\usepackage{xcolor}
\usepackage{array,tabularx,booktabs}
\usepackage[most]{tcolorbox}
\usepackage{tikz}

\setmonofont[Scale=0.88]{Menlo}
\setCJKmonofont{PingFang SC}
\xeCJKDeclareCharClass{CJK}{"201C, "201D, "2018, "2019}

\usepackage{enumitem}
\setlist{itemsep=2pt, topsep=4pt, parsep=0pt}

\definecolor{cmtgray}{RGB}{100,118,135}
\definecolor{rulegray}{RGB}{180,186,192}
\definecolor{introbg}{RGB}{252,248,238}
\definecolor{introline}{RGB}{206,162,90}
\definecolor{wechatgreen}{RGB}{7,193,96}
\definecolor{seriesnight}{RGB}{7,20,33}
\definecolor{seriesice}{RGB}{178,220,240}
\definecolor{seriescyan}{RGB}{46,163,214}

\newlength\ttcw
\AtBeginDocument{\settowidth\ttcw{{\footnotesize\ttfamily x}}}

\newtcolorbox{codebox}{
  breakable,enhanced,
  colback=black!3!white,frame hidden,
  borderline west={2.5pt}{0pt}{black!25},
  left=9pt,right=6pt,top=7pt,bottom=7pt,
  before skip=10pt,after skip=12pt,
  before upper={\setlength{\parskip}{0pt}\setlength{\parindent}{0pt}},
  fontupper=\footnotesize\ttfamily\frenchspacing\raggedright
}
\newcommand\codeline[1]{\par\noindent\hangindent=2.4em\hangafter=1 #1}
\newcommand\codeblank{\par\vspace{0.5\baselineskip}}

\newtcolorbox{noteblock}{
  breakable,enhanced,
  colback=blue!4!white,frame hidden,
  borderline west={2.5pt}{0pt}{blue!30!black!30},
  left=8pt,right=6pt,top=5pt,bottom=5pt,
  before skip=8pt,after skip=10pt,
  fontupper=\small
}

\newtcolorbox{introbox}{
  breakable,enhanced,
  colback=introbg,frame hidden,
  borderline west={2.5pt}{0pt}{introline},
  left=10pt,right=9pt,top=8pt,bottom=8pt,
  before skip=12pt,after skip=14pt,
  fontupper=\small\linespread{1.3}\selectfont
}

% 片段出处行：紧跟在代码框后，注明完整文件位置
\newcommand{\snipsource}[1]{%
  \par\vspace{-6pt}\noindent{\scriptsize\color{black!45}\(\triangleright\) 节选自仓库 \texttt{#1}，完整文件含更多小节与注释。}\par\vspace{4pt}}

% 章首插图
\newcommand{\chapterart}[2]{%
  \begin{center}
    \IfFileExists{figures/#1.png}{%
      \includegraphics[width=0.62\linewidth]{figures/#1.png}\par\vspace{0.4em}%
    }{%
      \begin{tcolorbox}[width=0.82\linewidth,height=4.2cm,colback=black!2,
        frame hidden,borderline={0.7pt}{0pt}{black!35,dashed},
        halign=center,valign=center]
        {\color{black!50}\small 〔插图位预留：运行 \texttt{gen\_figures.py} 生成〕}
      \end{tcolorbox}\par\vspace{0.4em}%
    }%
    {\small\color{black!60}\textbf{#2}}\par\vspace{0.5em}
    \begin{minipage}{0.82\linewidth}
      \scriptsize\color{black!45}\raggedright
      本图由 \texttt{gpt-image-2} 生成 · 提示词：\textit{\input{fragments/prompt-#1}}
    \end{minipage}
  \end{center}
  \vspace{0.4em}}

\usepackage[colorlinks,linkcolor=blue!35!black,urlcolor=blue!55!black]{hyperref}
\hypersetup{
  pdftitle={《工程本体论》·可扩展本体框架与 Agent Skill},
  pdfauthor={claude F5 max},
  bookmarksnumbered=true,
  unicode=true
}

\setlength{\parskip}{0.35em}
\usepackage{longtable}
